%%%%
\teteSndAtom

%%%% titre
\vspace*{-36pt}
\titreActivite{Le cortège électronique}

%%%% Objectifs
\begin{objectifs}
  \item Comprendre la structure du cortège électronique.
  \item Comprendre la règle de remplissage des couches électroniques.
\end{objectifs}

\begin{contexte}
  Un atome est constitué d'un noyau positif entouré d'électrons négatifs, avec autant d'électrons que de protons, l'atome étant neutre.
  
  \problematique{
    Comment les électron s'organisent autour du noyau ?
  }
\end{contexte}


%%%% Documents
\begin{doc}{Rangement des électrons}
  Quand on s'appelle hydrogène et qu'on a qu'un électron, pas besoin de ranger ses affaires.
  Mais quand on s'appelle uranium et qu'on en a 92 autour de soi, mieux vaut mettre un peu d'ordre dans ses électrons !
  
  \begin{wrapfigure}{r}{0.35\linewidth}
    \centering
    \vspace*{-24pt}
    \image{0.9}{images/atomes/schema_couches_oxygene}
    
    \legende{Schéma des couches et sous-couches électroniques de l'oxygène \isotope{O}{8}{} avec ses 8 électrons.
    Les sous-couches $s$ sont en traits pleins, les $p$ sont en pointillés.}
  \end{wrapfigure}
  
  C'est en 1913 que Bohr a l'idée de répartir les électrons d'un atome en différentes couches et sous-couches, en se basant sur les travaux de Planck.

  Les électrons s'organisent en \important{couches électroniques}, qui contiennent des \important{sous-couches électroniques.}
  
  Les couches électroniques sont numérotées \important{1, 2, 3.}
  Les sous couches sont repérées par des lettres : \important{s} ou \important{p}.
  Les sous-couches ne peuvent contenir qu'un nombre limité d'électrons.

  \begin{importants}  
    La \important{sous-couche s} ne peut contenir que \important{2 électrons} au maximum,
    alors que la \important{sous-couche p} ne peut contenir que \important{6 électrons} au maximum.
  \end{importants}
  
  La couche qui accueille les derniers électrons s'appelle \important{la couche externe}, les autres couches sont appelées les \important{couches internes}.
\end{doc}

\begin{doc}{Remplissage des couches électroniques}
  Attirés par les protons, les électrons vont remplir les sous-couches électroniques.
  
  \begin{importants}
    L'ordre de remplissage suit \important{la règle de Klechkowski :}
    \begin{center}
      \important{1s} \flecheLongue
      \important{2s} \flecheLongue \important{2p} \flecheLongue
      \important{3s} \flecheLongue \important{3p}
    \end{center}
  \end{importants}
  \begin{importants}  
    On appelle \important{configuration électronique} le remplissage des électrons dans chaque couches et sous-couches. Une sous-couche \important{doit être remplie} pour passer à la sous-couche suivante.
  \end{importants}

  \begin{importants}
    Pour décrire la configuration électronique, on note \important{chaque sous-couche qui contient des électrons} par son nom ($1s$, $2s$, $2p$, etc.), puis on note \important{le nombre d'électrons contenu dans la sous-couche en exposant}.
  
    \textit{Exemple :} la configuration électronique de l'atome d'oxygène \isotope{O}{8}{} est $1s^2 \; 2s^2 \; 2p^4$. 2 électrons dans la sous-couche 1s, 2 dans la sous-couche 2s, 4 dans la sous-couche 2p.
  \end{importants}
  

  \attention La somme des exposants doit être égale au nombre d'électrons de l'atome !
\end{doc}


%%%%
\newpage
\numeroQuestion
Compléter le tableau ci-dessous pour résumer l'occupation des différentes couches électroniques 

\vspace*{-12pt}
\begin{center}
  \begin{tblr}{
    colspec = {X[c] X[c] X[c] X[c] X[c] X[c]}, hlines, vlines,
    row{1} = {couleurSec-100, c}, column{1} = {l, couleurSec-50}
  }
    Couche & 1 & \SetCell[c=2]{c} 2 & & \SetCell[c=2]{c} 3 & \\
    Sous-couche & \correction{s} &
    \correction{s} & \correction{p} &
    \correction{s} & \correction{p} \\
    Nombre max. d'électron & \correction{2} &
    \correction{2} & \correction{6} &
    \correction{2} & \correction{6} \\
  \end{tblr}
\end{center}

\mesure
L'atome de silicium \chemfig{Si} possède $Z = 14$ protons.
Schématiser ci-dessous la répartition de ses électrons.

\begin{center}
  \pasCorrection{\image{0.4}{images/atomes/schema_couche}}
  \correction{\image{0.4}{images/atomes/schema_couche_silicium}}
\end{center}

\question{
  Donner la configuration électronique de l'atome de silicium.\label{exo:configuration}
}{
  On remplit les sous couches dans l'ordre,  Si : $1s^2\; 2s^2 2p^6\; 3s^2 3p^2$.
}[2]

\question{ 
  Indiquer, en justifiant, le nom de la couche externe de cet atome de silicium, ainsi que la ou les couches internes.\label{exo:couches}
}{
  La couche externe du silicium est la troisième couche, ses couches internes sont les couches 1 et 2.
}[3]

\question{
  Reprendre les questions~\ref{exo:configuration} et \ref{exo:couches} pour l'atome de Carbone \chemfig{C} ($Z = 6$). Quelles différences et ressemblances avec le silicium peut-on remarquer ?
}{
  C : $1s^2\; 2s^2 2p^4$.
  
  La couche externe du carbone est la couche 2, sa couche interne est la couche 1. Le silicium et le carbone ont les sous-couches remplies de la même façon sur leur couche externe.
}[5]
